\subsection{\textbf{Sprint 1}}

Ao final da Sprint 0, foi realizada uma reunião com o stakeholder, na qual ele aprovou a grande maioria das telas. Contudo, ele solicitou algumas alterações no fluxo e na tela inicial do aplicativo. Mais especificamente, foi solicitada a adição de dois novos botões na tela inicial: um que levaria à lista de frases e outro que permitiria a adição rápida de uma nova frase.

Na aula seguinte da apresentação, os squads foram divididos novamente para desenvolver as seguintes tarefas: o cadastro do usuário, a tela inicial e a refatoração das telas no Figma. Meu squad ficou responsável pela tela inicial. Junto com o \ac{ages} III Leonardo José Machado Canto, dividimos a tarefa da tela inicial em duas partes iniciais: os botões e o gráfico da meta diária. Cada dupla ficou responsável por uma dessas partes, e eu escolhi ser responsável pela criação dos demais componentes.

Como não tive tempo de estudar Flutter durante a Sprint 0, precisei focar exclusivamente nos estudos dirigidos durante a primeira semana da Sprint 1. Na primeira semana, junto com o Leonardo, criamos um esboço inicial, já que ainda estávamos aguardando o design do outro squad ser finalizado e aprovado pelo cliente. Infelizmente, nosso squad teve problemas de comunicação e falta de participação de alguns integrantes em suas tarefas, o que sobrecarregou nosso \textit{Squad Leader}, que acabou ajudando em várias partes da criação das telas.

Porém, a tela inicial só foi aprovada pelo stakeholder quatro dias antes do \textit{Freezing}. Logo no mesmo dia, comuniquei ao time que tentaria criar algo em Flutter parecido com o que o cliente havia aprovado, para que pudéssemos discutir possíveis alterações e melhorias na \ac{ui} do aplicativo. Acabei desenvolvendo parte do trabalho que, a princípio, não seria minha responsabilidade, mas se mostrou necessário. Neste dia, foram criados quase todos os componentes que compõem a tela inicial.

Durante a última aula antes da entrega, nos reunimos com os líderes do squad responsável pela refatoração das telas e definimos alguns pontos importantes. O principal é que o \textit{mockup} aprovado havia sido feito pelo próprio stakeholder, que demonstrou ser proativo, mas foi criado em um tamanho muito superior ao das outras telas, o que tornou inviável transpor fidedignamente essa tela para a aplicação.

Dois dias antes do \textit{Freezing}, ainda tínhamos pendências referentes à padronização de alguns elementos. Como havia sido solicitado pelo Leonardo, mas ninguém havia se mostrado interessado em pegar as tarefas, senti-me no dever de entregar os dois últimos itens que restavam para finalizar.

A demora na definição da tela acabou atrasando bastante a entrega do meu squad, e junto com outros fatores, como a falta de um ambiente de testes e um banco de dados para a retirada das informações do cadastro, nossa tarefa acabou por não ser aceita. Apesar disso, fico contente com o trabalho que realizei. Como mencionado anteriormente, eu não tinha conhecimento algum em Flutter e mesmo assim consegui estudar e contribuir muito durante a Sprint.